% !LW recipe = make-deck
% Level 6 (the tracks meet): lists represent tables. Theory: the table of
% a list, Boolean conditions as properties, the programs compute the
% operators, so the laws hold for lists. Lean: toTable, and the theorems.
% Build: make (rules in ../common/deck.mk).
\DocumentMetadata{
  lang        = en-GB,
  pdfstandard = ua-2,
  tagging     = on,
}
\documentclass{ltx-talk}

\input{preamble}

\title{Level 6: Lists Represent Tables}
\author{The Proofs and Programs Club}
\institute{Royal Holloway, University of London}

\begin{document}

\leveltitle{6}{Lists Represent Tables}{The two tracks meet}
  {After \levellink{5f} and \levellink{5t}}
  {Next: \levellink{7} or \levellink{9}}

\levelframe{6}{Level 6 at a Glance}{%
  \code{toTable}, the table a list represents.}{%
  that your list programs (programs track) compute the operators on
  tables (proofs track). So every law holds for your lists.}{%
  prove facts about programs by induction & \rds{fpil-tactics} \\
  find and use Lean's lemmas about lists, with \code{simp} &
    \rds{tpil-simp} \\
  split a goal on an \code{if} & \rds{tpil-split} \\
}

% ---------------------------------------------------------------- theory
\theorypart{Programs that compute tables}

\begin{frame}
  \frametitle{The Programs Compute the Operators}
  \defn{Definition 6.1}{The table of a list $l$ is $\mathit{table}(l) =
    \{t \mid t \text{ is in } l\}$: the one $l$ represents (Definition
    1.2). A condition $p$ with values in $\B$ gives the property ``$p(t)
    = \true$''.}
  \defn{Theorem 6.2}{For all lists $l$ and $m$:
    \[
      \begin{array}{@{}r@{\;=\;}l@{\qquad}r@{\;=\;}l@{}}
        \mathit{table}(\mathit{select}\ p\ l) & \sel{p}(\mathit{table}(l))
          & \mathit{table}(l \mathbin{+\!\!+} m) & \mathit{table}(l) \cup
          \mathit{table}(m) \\
        \mathit{table}(\mathit{project}\ f\ l) & f[\mathit{table}(l)] &
          \mathit{table}(\mathit{product}\ l\ m) & \mathit{table}(l)
          \times \mathit{table}(m) \\
      \end{array}
    \]}
  On the left, the programs of Levels 2F and 3F; on the right, the
  operators of Levels 2T to 4T.
  \theorytag{6}
\end{frame}

\begin{frame}
  \frametitle{One Induction, and Every Law Comes for Free}
  \emph{Proof of Theorem 6.2, for $\mathit{select}$.} By induction on $l$.
  For $[\,]$, both sides are empty. For $t :: l$, if $p(t) = \true$, the
  table of $t :: \mathit{select}\ p\ l$ is $\{t\} \cup
  \sel{p}(\mathit{table}(l))$, by the assumption, which is
  $\sel{p}(\{t\} \cup \mathit{table}(l))$. If $p(t) = \false$, $t$ is
  dropped on both sides. \proofend
  \defn{Corollary 6.3}{Every law of Levels 2T to 4T holds for the tables
    of lists. For example: $\mathit{table}(\mathit{select}\ p\
    (\mathit{select}\ q\ l)) = \mathit{table}(\mathit{select}\ (p \wedge
    q)\ l)$.}
  \emph{Proof.} By Theorem 6.2 twice, both sides are
  $\sel{p}(\sel{q}(\mathit{table}(l)))$ and $\sel{p \wedge
  q}(\mathit{table}(l))$, equal by law 1. \proofend
  \theorytag{6}
\end{frame}

% ---------------------------------------------------------------- Lean
\leanpart{The two tracks, joined}

\begin{frame*}
  \frametitle{One Line Joins the Two Tracks}
  \leanblock{totable.lean}
  \begin{itemize}
    \item One file now holds both tracks: your programs (\code{select},
      \code{union}, \ldots) and \code{Sets} (\code{Sets.select}, \ldots).
    \item \code{simp} knows that \code{t ∈ l ++ m} means \code{t ∈ l ∨ t ∈
      m}, and finishes.
  \end{itemize}
  \slideref{\lnk{tpil-simp}}
  \leantag{6}
\end{frame*}

\begin{frame*}
  \frametitle{Your Turn: select Computes a Selection}
  \leanblock{represent.lean}
  \begin{itemize}
    \item \ex{6.1} Induction on \code{l}; in the \code{cons} case,
      \code{cases h : p row}.
    \item Unfold with \code{simp [toTable, select, Sets.select, h] at ih
      ⊢}, then finish with \code{ih} and logic (Level 2T).
  \end{itemize}
  \slideref{\lnk{fpil-tactics} \rfsep \lnk{tpil-split}}
  \leantag{6}
\end{frame*}

\begin{frame*}
  \frametitle{Your Turn: project Computes an Image}
  \leanblock{project.lean}
  \begin{itemize}
    \item \ex{6.2} The same, for \code{project}: the \code{∃} of Level 4T
      comes back.
    \item \ex{6.3} Stretch: \code{product}; and \code{diff}, assuming
      \code{[LawfulBEq α]}, so that \code{==} is equality.
  \end{itemize}
  \slideref{\lnk{tpil-exists}}
  \leantag{6}
\end{frame*}

\begin{frame*}
  \frametitle{Your Turn: A Law for Lists, with No New Induction}
  \leanblock{law.lean}
  \begin{itemize}
    \item \ex{6.4} Corollary 6.3. Start with \code{simp only
      [toTable\_select, Sets.select]}, then finish with logic, or with
      your \code{Sets.select\_select}.
    \item \ex{6.5} Pick another law of Levels 2T to 4T, and prove it for
      lists the same way.
  \end{itemize}
  \leantag{6}
\end{frame*}

\end{document}
